\FloatBarrier
\section{Conclus\~ao}

Este estudo tratou a convers\~ao de pontos nucleares em caixas como vari\'avel
cient\'ifica. Para a RQ1, observou-se um plat\^o de F1 entre 96 e 192 pixels, com
compromissos distintos de AP50, revoca\c{c}\~ao e contagem. Para a RQ2, YOLO11s
foi o melhor ponto operacional, embora a vantagem sobre YOLO11n tenha sido
pequena e YOLO11m tenha apresentado maior revoca\c{c}\~ao. Para a RQ3, o modelo
selecionado obteve AP50 0,8723, F1 0,8190 e MAE 4,68 no teste retido; o erro
absoluto aumentou em imagens mais densas.

A conclus\~ao principal n\~ao \'e que 144 pixels ou YOLO11s sejam universais. O
resultado relevante \'e que a defini\c{c}\~ao espacial criada a partir de pontos
altera de forma mensur\'avel detec\c{c}\~ao e contagem e, portanto, precisa fazer
parte do protocolo reportado. O pipeline produzido oferece uma refer\^encia
reprodut\'ivel para esse tipo de auditoria na CRIC. Sua extens\~ao natural \'e
substituir a geometria fixa por supervis\~ao diretamente pontual e comparar com
uma amostra de contornos anotados.

\section*{Artefatos de Reprodu\c{c}\~ao}

Os CSVs usados pelos gr\'aficos, configura\c{c}\~oes de treinamento, checkpoints
selecionados e rotinas de avalia\c{c}\~ao est\~ao dispon\'iveis no reposit\'orio
p\'ublico do projeto. O reposit\'orio inclui as parti\c{c}\~oes por imagem (semente
42), os metadados das pseudo-caixas para cada valor de $s$ e os arquivos de
predi\c{c}\~ao do conjunto de teste, permitindo reproduzir as m\'etricas reportadas
sem re-treinamento.
